\section{The coupled elliptic construction and its limiting geometry}
\label{pde:section}

We now solve the system of Section~\ref{alg:section}.  Its preceding
estimates control the unknowns and their first derivatives, but the trace
operator has a coefficient whose modulus of continuity is not initially
controlled.  We prove the particular gradient estimate needed for that
operator.  This estimate, a scalar potential argument for the second
equation, and continuation give smooth solutions.  We then establish the
end behavior, mass cost, and height barriers that keep the region requiring
charged curvature control on the original exterior.  No charged
solvability theorem is used.

The prepared data, dimension, and \(\varepsilon>0\) are fixed first, and
\(N\) is subsequently chosen sufficiently large.  We use
\eqref{alg:definitions} and \eqref{alg:system}, including
\[
 k=n-1,\qquad \ell=e^{-N\varepsilon},\qquad
 \eta=\ell^{3/2},\qquad h=\eta f.
\]
The notation \(\mathcal P_N\) denotes a possibly changing polynomial in
\(1+N\), with coefficients depending on the fixed data, dimension, and
\(\varepsilon\).  Local elliptic constants may depend arbitrarily on fixed
\(N\).  They are not used as uniform constants in the limit \(N\to\infty\).

\subsection{A gradient estimate without a coefficient modulus}

\begin{lemma}[The special scalar gradient estimate]
\label{pde:special-gradient}
Let \(g\) be a smooth uniformly positive metric on a coordinate ball, or
on a smooth coordinate half-ball with a distinguished boundary face,
with bounded smooth coordinate geometry.  Suppose an individually smooth
function \(z\) satisfies
\begin{equation}
 \left(g^{-1}-(1-\chi)e_z\otimes e_z\right):\nabla^2 z=G,
 \qquad e_z=\frac{\nabla z}{|\nabla z|_g},
 \qquad c\le\chi\le1,
 \label{pde:special-equation}
\end{equation}
where \(c>0\), \(|\nabla z|_g\le M\), and \(|G|\le C\).
The axis at a gradient zero may be chosen arbitrarily.  The coefficient
\(\chi\) may be measurable, with no assumed modulus of continuity.
In the half-ball case assume a constant value of \(z\) on the
distinguished face.

There are \(\alpha\in(0,1)\) and uniform local \(C^\alpha\) bounds for
\(\nabla z\), in the interior or up to that face.  On smaller balls or
half-balls one also has
\begin{equation}
 \int_{B_r}|\nabla^2z|_g^2\,\dd V_g
 \le C_1r^{n-2+2\alpha},\qquad 0<r<r_1.
 \label{pde:hessian-morrey}
\end{equation}
Constants depend on \(n,c,M,C\), the coordinate geometry and localization,
but not on the modulus of \(\chi\) or higher derivatives of the individual
solution.
\end{lemma}

\begin{proof}
We give the boundary reduction and both improvement arguments.

\paragraph{A local Hessian estimate.}
Normalize the metric to the identity at a ball center and restrict to a
sufficiently small ball.  The coordinate principal matrix \(a\) then
satisfies
\[
 |I-a|_{\mathrm F}\le1-c/2,
\]
where the norm is Frobenius; connection terms are lower-order terms.
The full-space operator \(D^2\Delta^{-1}\), with the Frobenius Hessian
norm, has \(L^2\) norm one by Fourier transformation.  Its boundedness on
nearby \(L^s\) spaces and interpolation give \(s_0>2\) for which its
norm times \(1-c/2\) is less than one.  Absorbing the perturbation from
the Laplacian and localizing therefore gives, for \(s=2,s_0\),
\begin{equation}
 \|D^2Y\|_{L^s(B_r)}
 \le C_s\left(\|a:D^2Y\|_{L^s(B_{2r})}
                 +r^{-2}\|Y\|_{L^s(B_{2r})}\right).
 \label{pde:cordes-local}
\end{equation}
The cutoff commutator contains \(r^{-1}DY\) and \(r^{-2}Y\).
The interpolation estimate
\(\|DY\|_s\le\delta r\|D^2Y\|_s+C_\delta r^{-1}\|Y\|_s\),
and absorption on nested balls, remove the former.  This argument uses
only measurable principal coefficients.  The Laplacian estimates are the
classical singular-integral bounds
\cite{CalderonZygmund1952,GilbargTrudinger2001}.

\paragraph{Reflection and its flux correction.}
Subtract the constant face value.  In Gaussian coordinates write
\[
 g=\dd t^2+h_{ab}(y,t)\,\dd y^a\dd y^b,\qquad t\ge0.
\]
Extend \(z\) oddly, \(h\) and \(\chi\) evenly, and \(G\) oddly across
\(t=0\).  For \(R=\operatorname{diag}(I,-1)\), the reflected gradient
and Hessian are \(-R\nabla z\) and \(-R^T(\nabla^2z)R\), while the
principal matrix is \(RaR\).  Thus the equation holds almost everywhere.
Tangential derivatives of the face value vanish, so the extension is
individually \(C^1\) and \(W^{2,s}\) for finite \(s\).  The reflected
metric is Lipschitz and smooth on each side, sufficient for
\eqref{pde:cordes-local}.

There is one boundary contribution when differentiating the gradient norm.
Put
\[
 p_z=\nabla z,\quad U_z=|p_z|_g^2/2,\quad
 A=g^{-1}-(1-\chi)e_z\otimes e_z,\quad
 \mathcal F=(\nabla^2z-(\Delta_gz)g)\nabla z.
\]
The equation and the contracted Hessian identity give
\begin{align}
 A\,\dd U_z-Gp_z&=\mathcal F,
 \label{pde:gradient-flux}\\
 \operatorname{div}_g\mathcal F
 &=|\nabla^2z|_g^2-(\Delta_gz)^2+\operatorname{Ric}_g(p_z,p_z)
 \label{pde:bochner-flux}
\end{align}
on each smooth side.  Neither formula differentiates \(\chi\).
Let \(q=z_t(y,0)\),
\(\mathrm{II}_{ab}=\tfrac12\partial_t h_{ab}(y,0)\),
\(H=\operatorname{tr}_h\mathrm{II}\), and \(J=\sqrt{\det h}\).
Constant face data give \((\nabla^2z)_{ab}=q\mathrm{II}_{ab}\).
Thus the normal traces of \(\mathcal F\) are \(-Hq^2\) and \(+Hq^2\).
The coordinate divergence of \(J\mathcal F\) has the plane density
\(-2J(y,0)Hq^2\delta_{\{t=0\}}\).  It is canceled by
\begin{equation}
 \mathcal C
 =2J(y,0)H(y)q(y)^2\,\mathbf1_{\{t>0\}}\partial_t.
 \label{pde:face-correction}
\end{equation}
This vector has only a normal component.  Its divergence differentiates
only in \(t\), and requires no tangential derivative of \(q\).
For each localized ball, extend the face density
\(2J(y,0)H(y)q(y)^2\) by zero in the tangential variables outside
the portion of the face lying in that ball.  Since \(\mathcal C\)
has only a normal component, this localization creates no additional
distributional divergence and its bound uses only the gradient bound
on that face portion.

Normalize \(|p_z|_g\le1\) and set \(s=1-|p_z|_g^2\).  The corrected flux is
\[
 JA\,Ds+B=-2(J\mathcal F+\mathcal C),\qquad
 B=2JGp_z-4J(y,0)Hq^2\mathbf1_{\{t>0\}}\partial_t.
\]
Convexity gives
\[
 (\Delta_gz)^2
 =\big((1-\chi)\nabla^2z(e_z,e_z)+G\big)^2
 \le(1-\chi)|\nabla^2z|_g^2+G^2/\chi.
\]
Consequently, in distributions,
\begin{equation}
 \operatorname{div}(JA\,Ds+B)
 \le-2cJ|\nabla^2z|_g^2
       +2J\big(G^2/c+|\operatorname{Ric}_g||p_z|_g^2\big).
 \label{pde:deficit-inequality}
\end{equation}
In an interior chart take \(\mathcal C=0\).
Radius-\(r\) rescaling multiplies face curvature by \(r\), smooth-side
curvature by \(r^2\), and the forcing by \(r/M_1\) if gradients are
divided by \(M_1\).  Thus the additional bounded flux is small under
the stated small-error rescaling.  The construction also applies to
balls centered near the reflected face.

\paragraph{Gradient drop or affine entry.}
Fix \(r_*\in(0,1/32)\) and \(b_0>0\).  There exist
\(k_*\in(r_*,1)\) and a small error threshold such that a normalized
solution on \(B_1\), with \(|\nabla z|_g\le1\), satisfies either
\begin{align}
 &\sup_{B_{r_*}}|\nabla z|_g\le k_*,
 \label{pde:gradient-drop}\\
 &\|z-L\|_{L^\infty(B_{r_*})}\le b_0r_*,
 \qquad \tfrac12\le|DL|\le2
 \label{pde:affine-entry}
\end{align}
for some affine \(L\).  Small errors comprise forcing, metric and
connection errors, smooth-side curvature, and reflected-face curvature.
Both \(k_*\) and the threshold may depend on \(b_0\).

Otherwise choose errors tending to zero, with the gradient supremum on
\(B_{r_*}\) tending to one and without affine entry.  Normalize
\(z_j(0)=0\).  Equation~\eqref{pde:cordes-local} gives uniform local
\(W^{2,2}\) bounds and hence uniform local \(L^2\) bounds on
\(Ds_j\), where \(s_j=1-|\nabla z_j|_{g_j}^2\).
Equation~\eqref{pde:deficit-inequality} reads
\[
 \operatorname{div}(a_jDs_j+B_j)
 \le-\kappa|D^2z_j|^2+\epsilon_j,\qquad
 \|B_j\|_\infty+\epsilon_j\longrightarrow0
\]
with uniform ellipticity and \(\kappa>0\).
Weak Harnack, after discarding the Hessian term, gives
\(s_j\to0\) in measure on \(B_{1/2}\), since
\(\inf_{B_{r_*}}s_j\to0\).

For completeness, convergence of \(Ds_j\) needs a separate test.
With a cutoff \(\zeta\) supported in \(B_{1/2}\), test the inequality
by \(\zeta^2(b-s_j)_+\), \(0<b<1\).  Ellipticity and Young give
\begin{align}
 \int_{\{s_j<b\}}\zeta^2|Ds_j|^2
 \le C\bigg(&b^2\int|D\zeta|^2+\int\zeta^2|B_j|^2
 \nonumber\\
 &+b\int\zeta|D\zeta||B_j|
   +\epsilon_j b\int\zeta^2\bigg).
 \label{pde:lower-truncation}
\end{align}
Its limsup on a smaller compact ball is at most \(Cb^2\).
On \(\{s_j\ge b\}\), measure convergence and the uniform \(L^2\)
gradient bound make the \(L^1\) norm of \(Ds_j\) tend to zero.
On \(\{s_j<b\}\), the limsup is at most \(Cb\).
Sending \(b\) to zero gives local \(L^1\) convergence of \(Ds_j\).
A fixed cutoff in the full deficit inequality now yields
\(D^2z_j\to0\) in \(L^2(B_{1/3})\).
Poincar\'e makes the gradients converge to a constant vector, of length
one because \(s_j\to0\).  Uniform Lipschitz control gives uniform
convergence on smaller balls to the corresponding affine function,
a contradiction.  Here and below we use the usual inhomogeneous weak
Harnack and scalar divergence estimates
\cite{GilbargTrudinger2001}.

\paragraph{Improvement near a nonzero affine function.}
There are \(\alpha_2>0\), \(\lambda\in(0,1/8)\), and sufficiently
small \(b_0,\delta'>0\) with this property: if
\[
 \|z-L\|_{L^\infty(B_1)}\le b\le b_0,\qquad
 \tfrac14\le|DL|\le4,
\]
the gradient is uniformly bounded, and metric, connection and forcing
errors are at most \(b\delta'\), then an affine \(L'\) satisfies
\begin{equation}
 \|z-L'\|_{L^\infty(B_\lambda)}
 \le b\lambda^{1+\alpha_2},\qquad |DL'-DL|\le Cb.
 \label{pde:affine-improvement}
\end{equation}
Indeed \(Y=(z-L)/b\) has bounded supremum and bounded local
\(W^{2,s_0}\) norm by \eqref{pde:cordes-local}.  The normalized
connection term is small because it uses the original gradient bound
before division by \(b\).  As \(b_0,\delta'\to0\),
\(Dz-DL=bDY\to0\) locally in measure.  Freeze only the axis at
\(e_0=DL/|DL|\), leaving \(\chi\) measurable.  Bounded coefficient
differences converging in measure, and the \(L^{s_0}\) Hessian bound,
give
\begin{equation}
 \|\Delta_{e_0^\perp}Y+\chi\partial_{e_0}^2Y\|_{L^2(B_{3/4})}
 \longrightarrow0.
 \label{pde:axis-residual}
\end{equation}
The Dirichlet Laplace inverse on the convex ball \(B_{3/4}\) has
\(\|D^2v\|_2\le\|\Delta v\|_2\), by integration by parts.
Thus the frozen operator has a zero-Dirichlet inverse in \(W^{2,2}\)
by a contraction of norm at most \(1-c\).
Subtract the resulting small correction to obtain \(Y_0\), uniformly
bounded in \(W^{2,2}\), with
\[
 \Delta_{e_0^\perp}Y_0+\chi\partial_{e_0}^2Y_0=0.
\]

The derivative \(\zeta_0=\partial_{e_0}Y_0\in W^{1,2}\) satisfies
\begin{equation}
 \operatorname{div}\left[
 (I-e_0\otimes e_0+\chi e_0\otimes e_0)D\zeta_0\right]=0
 \label{pde:frozen-derivative}
\end{equation}
in distributions.  Differentiation acts on the \(L^2\) flux
\(\chi\partial_{e_0}\zeta_0\); no classical derivative of \(\chi\)
is being asserted.  Uniform scalar divergence estimates give
\(\zeta_0\in C^{\alpha_1}\) for some \(\alpha_1\in(0,1)\).
Caccioppoli after subtraction of a local value gives
\[
 \int_{B_r}|D\zeta_0|^2\le Cr^{n-2+2\alpha_1}.
\]
The source in
\(\Delta Y_0=(1-\chi)\partial_{e_0}\zeta_0\)
therefore has \(L^1\) mass at most \(Cr^{n-1+\alpha_1}\).
The gradient of its localized Newton potential is \(C^{\alpha_1}\):
for points separated by \(a\), near annuli of radius \(R\le a\)
contribute \(CR^{\alpha_1}\), and the far-field gradient difference
contributes \(CaR^{\alpha_1-1}\) for \(R\ge a\).
Both dyadic sums converge.  Subtracting this potential leaves a harmonic
function on a smaller ball.  All first derivatives of \(Y_0\) thus
have a uniform H\"older bound.

The unfrozen \(Y\) has a uniform interior H\"older modulus by the scalar
Krylov--Safonov estimate with bounded forcing
\cite[Section~9.8]{GilbargTrudinger2001}; see also \cite{Mooney2019}
for the homogeneous mechanism.  The \(L^2\)-small correction
\(Y-Y_0\) is consequently small uniformly on a smaller ball, using
the uniform moduli of both terms.  Taylor approximation of \(Y_0\),
first choosing \(0<\alpha_2<\alpha_1\) and small \(\lambda\), and
then \(b_0,\delta'\), proves \eqref{pde:affine-improvement}.

\paragraph{Iteration and Hessian growth.}
At each center begin with a fixed gradient normalization and sufficiently
small radius for both error thresholds.  On each gradient drop reduce
the radius by \(r_*\) and the normalization by \(k_*\).
Since \(r_*/k_*<1\), normalized forcing stays small; geometric errors
also decrease.  If affine entry never occurs, the approximation by a
constant has order \(r^{1+\alpha_d}\), where
\(\alpha_d=\log k_*/\log r_*>0\).

After affine entry rescale its smaller ball to unit size and keep the
same gradient normalization.  Iterate \eqref{pde:affine-improvement}
at radii \(\lambda^i\), with errors
\(b_i=b_0\lambda^{i\alpha_2}\).
Only constants, not affine functions, are subtracted when rescaling
the structural equation, so its axis remains that of the actual
gradient.  Forcing and geometric errors decrease like \(\lambda^i\),
faster than \(b_i\).  The slope increments \(Cb_i\) are summable;
small \(b_0\) keeps slopes in \([1/4,4]\).
Combining the two regimes gives uniform affine approximation of order
\(r^{1+\alpha}\), also across the entry scale, for
\(\alpha\le\min(\alpha_d,\alpha_2)\).
The affine Campanato criterion gives the gradient H\"older estimate.
The same iteration applies in corrected reflected charts.

Finally subtract the first affine jet in \eqref{pde:cordes-local}.
Its remainder is \(O(r^{1+\alpha})\), and the forcing and connection
terms are bounded.  The squared \(s=2\) estimate gives
\[
 \int_{B_r}|D^2z|^2\le C(r^n+r^{n-2+2\alpha}),
\]
which implies \eqref{pde:hessian-morrey}, with equivalent covariant
norms and the same argument on half-balls.
\end{proof}

\subsection{The second equation and coupled regularity}

\begin{lemma}[A bounded solution with natural gradient growth]
\label{pde:natural-growth}
On a ball suppose a bounded, individually smooth function \(Z\) solves
\[
 -\operatorname{div}(A\,DZ)=\operatorname{div}B+F,
\]
where \(A\) is symmetric, bounded and uniformly elliptic, \(B\) is
bounded, and
\[
 |F|\le a_0|DZ|^2+\mu,\qquad
 \mu(B_r(x))\le C_0 r^{n-2+\alpha_0}
\]
for some \(\alpha_0\in(0,1)\), uniformly at the relevant centers.
Here \(\mu\) can be a nonnegative measure; in the application its
density is individually bounded.  Then \(Z\) has a uniform local
H\"older modulus.  The same conclusion holds up to a constant
Dirichlet face if the equation can be reflected with bounded
coefficients, bounded divergence flux, and the same source growth.
\end{lemma}

\begin{proof}
Choose a fixed large \(L\) and set \(Q_\pm=e^{\pm LZ}\).
The distributional chain rule gives
\begin{align*}
 -\operatorname{div}(A\,DQ_\pm)
 ={}&\operatorname{div}(\pm Le^{\pm LZ}B)\\
 &+e^{\pm LZ}\big(\pm LF-L^2 A\,DZ\cdot DZ-L^2B\cdot DZ\big).
\end{align*}
If the lower ellipticity bound is \(\lambda_0\), choose
\(L>2a_0/\lambda_0\).  Young's inequality absorbs the bounded-flux
cross term into the remaining negative gradient square.
The boundedness of \(Z\) gives
\begin{equation}
 -\operatorname{div}(A\,DQ_\pm)
 \le\operatorname{div}B_\pm+C(\dd x+\mu),
 \qquad \|B_\pm\|_\infty\le C.
 \label{pde:exponential-subsolutions}
\end{equation}
The matrix on the left is the same for both signs.

On \(B_r\), let \(w_\pm\) be the zero-Dirichlet potential of the
right side of \eqref{pde:exponential-subsolutions}.  Scalar Dirichlet
Green estimates for bounded measurable uniformly elliptic coefficients
give
\begin{equation}
 \|w_\pm\|_\infty\le C(r+r^{\alpha_0}).
 \label{pde:green-potential}
\end{equation}
Here is the dependence on the source hypotheses.  The bound
\(\mathcal G(x,y)\le C|x-y|^{2-n}\) and dyadic annuli give
\(Cr^{\alpha_0}\) for the measure source.  The exterior-pole gradient
energy bound
\(\int_{\{|x-y|>\rho\}}|D_y\mathcal G(x,y)|^2
 \le C\rho^{2-n}\)
gives \(L^1\) gradient mass \(C\rho\) on an annulus of scale \(\rho\).
Summing gives \(Cr\) for bounded divergence data.
These scalar Green bounds hold in every \(n\ge3\);
see \cite[Corollary~4.1 and Theorem~3.6]{KangKim2010}.
One can first use the individually bounded densities here and then
pass to their uniform growth bounds.

Put \(\epsilon_r=C(r+r^{\alpha_0})\), and define
\[
 U_\pm=\sup_{B_r}Q_\pm+\epsilon_r-Q_\pm+w_\pm.
\]
They are nonnegative supersolutions.  Let \(M=\sup_{B_r}Z\) and
\(m=\inf_{B_r}Z\).  On at least half of an intermediate ball either
\(Z\le(M+m)/2\) or \(Z\ge(M+m)/2\).
In the first case the uncorrected plus deficit is bounded below by
a fixed multiple of \(M-m\); in the second case the minus deficit is.
The multiplicative constant is positive because \(Z\) is bounded.
Weak Harnack for the corresponding \(U_\pm\) therefore lowers the
supremum, or raises the infimum, on \(B_{r/4}\), up to \(C\epsilon_r\).
Thus
\[
 \operatorname{osc}_{B_{r/4}}Z
 \le(1-c_0)\operatorname{osc}_{B_r}Z+C(r+r^{\alpha_0}).
\]
Iteration gives a H\"older modulus.  The same proof applies after the
stated reflection.
\end{proof}

\begin{lemma}[Estimates for smooth fixed points above the floor]
\label{pde:coupled-estimates}
Fix \(N\) and a sufficiently distant truncation of \(\mathcal M_N\).
Every smooth solution of the system \eqref{alg:system}, or of its
specified homotopy, satisfying \(t>-\varepsilon\), has local smooth
bounds depending on fixed \(N\) and the background geometry.
The bounds hold up to the outer zero-Dirichlet face and are uniform
as that face tends to infinity, on fixed compact sets and uniformly
sized end patches.
\end{lemma}

\begin{proof}
Lemma~\ref{alg:floor} and Proposition~\ref{alg:distortion-proposition} give
\[
 -\varepsilon<t\le Z\le\mathcal P_N,\qquad
 \ell\le l\le e,\qquad |f|+|Df|\le e^{\mathcal P_N}.
\]
At fixed \(N\), the eigenvalue
\(\chi=kd/(k+pv)\) is consequently bounded away from zero.
The normalized trace equation is
\begin{equation}
 A_\chi:\nabla^2f
 =\frac{\sqrt D}{l}
       \big(\eta f-\operatorname{tr}_{A_\chi}K\big).
 \label{pde:normalized-trace}
\end{equation}
Its right side is bounded.  Lemma~\ref{pde:special-gradient}
therefore bounds \(Df\) in \(C^\alpha\), including at the outer face,
and gives \eqref{pde:hessian-morrey} for \(f\).

Write the other equation in the form of
Lemma~\ref{pde:natural-growth}, with
\[
 A=kuA_\chi,\qquad B=uA_\chi K(w,\cdot).
\]
The sign of the scalar source is immaterial for its absolute bound.
The implicit relation for \(t\) gives
\[
 |Dt|\le C_N(1+|DZ|+|D^2f|).
\]
The explicit quadratic formula for \(\mathcal T\) then bounds the
absolute scalar source by
\[
 C_N(1+|DZ|^2)+C_N|D^2f|^2.
\]
All lower-order fluxes are bounded.  The Hessian estimate supplies
\[
 \int_{B_r}|D^2f|^2\le C_Nr^{n-2+2\alpha};
\]
we may decrease the exponent when applying
Lemma~\ref{pde:natural-growth}.  This proves a uniform H\"older modulus
for \(Z\).

At the outer face extend \(Z\) oddly and the divergence matrix by
\(A^-=RA^+R\).  Extend the total flux by \(-R\), and the scalar
source oddly.  The normal flux is continuous, so its distributional
divergence has no plane mass.  Boundedness and the source-growth
estimates are preserved.  This justifies the use of the boundary
case of Lemma~\ref{pde:natural-growth}.

The coefficients and right side of
\eqref{pde:normalized-trace} are now H\"older, since they are smooth
functions of \(Z,Df\) on the bounded parameter range.  Interior and
Dirichlet Schauder estimates give \(f\) through second derivatives.
Expanding the divergence equation for \(Z\), its principal matrix is
H\"older and its remaining source is bounded by \(C_N(1+|DZ|^2)\).
For finite \(s>n\), use the local interpolation inequality
\begin{equation}
 \|DZ\|_{L^{2s}(Q)}^2
 \le C\,\operatorname{osc}_{Q'}Z\,\|D^2Z\|_{L^s(Q')}
       +C_{Q,Q'}(\operatorname{osc}_{Q'}Z)^2
 \label{pde:quadratic-interpolation}
\end{equation}
on concentric ball or half-ball patches \(Q\Subset Q'\).
It follows by scaled derivative interpolation after subtraction of
a constant.  The already proved H\"older estimate makes its Hessian
coefficient uniformly small on sufficiently small fixed patches.
Local nondivergence \(W^{2,s}\) estimates therefore absorb the
quadratic term.  The leading estimate constants stay bounded as
patches shrink, since the principal matrices have a uniform H\"older
modulus.  Lower-order constants may depend on the selected scale.
Use bounded overlap of patch enlargements and the supremum of their
local Hessian norms to carry out the absorption; these norms are
individually finite on every truncation.  This proves \(W^{2,s}\)
bounds for \(Z\).

Schauder estimates and differentiation now bootstrap both functions.
At each differentiation the \(f\) equation is solved first with the
known regularity of \(Z\), and then the expanded \(Z\) equation gains
the next derivative.  These are the usual local and Dirichlet
elliptic estimates \cite{GilbargTrudinger2001}.  All charts, including
uniformly sized end charts and outer boundary charts, have controlled
geometry.  Hence the constants on fixed compact sets do not depend
on the truncation radius.
\end{proof}

\subsection{Solving the trace equation and the coupled system}

\begin{lemma}[The trace equation for an arbitrary trial scalar]
\label{pde:trial-trace}
On a fixed smooth truncation, fix \(Z\in C^{1,\alpha}\) with zero
outer trace.  The trace equation \(F=\eta f\), with zero outer
Dirichlet value for \(f\), has a unique solution in \(C^{3,\alpha}\).
The solution depends continuously on the trial \(Z\) and on the
tensor homotopy.  On bounded sets of trials these solutions have
bounded \(C^{3,\alpha}\) norms.  No floor condition on the trial
is needed.
\end{lemma}

\begin{proof}
For any prescribed \(Z,Df\), the defining relation for \(t\) is strictly
increasing in \(t\), so the coefficients are smooth functions of these
variables.  Continue the normalized equation
\eqref{pde:normalized-trace} from \(\Delta f=f\), by straight
interpolation of the two operators.  Its height coefficient is strictly
negative when all terms are placed on the left.  At a positive maximum
the gradient vanishes, \(A_\chi=I\), and the maximum test gives
\(f\le\|\operatorname{tr}K\|_\infty/\eta\); the analogous bound holds
at a minimum.  The same height bound works throughout the interpolation.

We supply the gradient bound, since ellipticity must not be assumed
before it.  Put \(\sigma=|Df|\).  For bounded trial \(Z\), as
\(\sigma\to\infty\), the implicit relation has \(t<0\) and hence
\(l=e^{Nt}\), so
\[
 t=-\frac{\log\sigma}{k+N}+O(1),\qquad
 d\asymp\sigma^{-2k/(k+N)}.
\]
For \(N>k\) this gives
\begin{equation}
 \chi\ge\frac{c_N}{1+\sigma},\qquad
 \frac{\sqrt D}{l}\le C_N(1+\sigma).
 \label{pde:trial-growth}
\end{equation}
The constants are uniform on a bounded set of trials.  Compact
\(\sigma\) ranges are included by positivity and continuity.
The normalized right side, using the height bound, is at most
\(C_N(1+\sigma)\) in absolute value.  Straight operator interpolation
has the same properties: its axial eigenvalue is at least \(\chi\).

Choose a short distance interval \([0,s_0]\) and a modulus
\[
 \psi(0)=0,\qquad \psi''=-K_0(\psi')^2,\qquad \psi'>0.
\]
Choose \(K_0\) large first, then \(s_0\) sufficiently small for the
distance geometry and so that \(\psi'\ge1\) can hold throughout.
Finally choose the initial slope sufficiently large that
\(\psi(s_0)\) exceeds twice the height bound.  This is possible from
\[
 \psi(s)=K_0^{-1}\log(1+K_0\psi'(0)s).
\]
On a boundary collar, \(\pm\psi(\operatorname{dist}(\cdot,\partial))\)
are upper and lower barriers.  The negative axial term
\(\chi\psi''\), by \eqref{pde:trial-growth}, has size at least a
fixed multiple of \(K_0\psi'\), and dominates the linear-growth
right side and the bounded second derivatives of the distance.

For the interior estimate consider a positive maximum of
\[
 f(x)-f(y)-\psi(\operatorname{dist}(x,y))
\]
with distance less than \(s_0\), using the smooth ambient metric
extension across the outer boundary.  Contacts at a boundary endpoint
are excluded by the preceding barriers; contacts at distance \(s_0\)
are excluded by the height bound.  At an interior contact the two
gradients are parallel along the short minimizing geodesic and have
length \(\sigma=\psi'\).  Parallel transverse endpoint variations
bound the difference of the transverse Hessian traces by
\(C\operatorname{dist}(x,y)\sigma\).  The separate axial variations
give
\[
 \nabla^2f_x(e,e)\le\psi'',\qquad
 -\nabla^2f_y(e,e)\le\psi''.
\]
Thus the difference of the two operators is at most
\[
 C s_0\sigma+(\chi_x+\chi_y)\psi''
 \le C s_0\sigma-c_NK_0\sigma.
\]
For large \(K_0\) this contradicts the absolute linear-growth bound
for the difference of their right sides.  The same argument with
\(x,y\) exchanged proves the Lipschitz bound.  It is valid throughout
the operator interpolation.

The gradient bound gives uniform ellipticity.  Apply
Lemma~\ref{pde:special-gradient}, then Schauder to obtain \(C^{2,\beta}\)
for some \(\beta>0\).  These bounds make \(Df\) Lipschitz, so the
original \(C^{1,\alpha}\) trial and the smooth coefficient functions
give \(C^\alpha\) coefficients; Schauder then yields \(C^{2,\alpha}\)
and, by one differentiation, \(C^{3,\alpha}\).
The Dirichlet linearization has a strictly negative zero-order term.
The maximum principle and the usual elliptic Fredholm alternative
therefore make it invertible.  The same maximum principle, applied
to the linear equation for the difference, gives uniqueness.
Continuation is open by this invertibility and closed by the uniform
estimates.  The implicit-function theorem and uniqueness also give
continuous dependence in the stated spaces.
\end{proof}

\begin{proposition}[A filled solution]
\label{pde:filled-solve}
For fixed prepared data and \(\varepsilon>0\), all sufficiently large
\(N\) admit a smooth solution of \eqref{alg:system} on the entire
one-ended filled manifold \(\mathcal M_N\), with
\[
 t>-\varepsilon,\qquad
 t\le Z\le\mathcal P_N,\qquad
 \ell\le l\le e,\qquad
 |f|+|\nabla f|\le e^{\mathcal P_N}.
\]
The functions \(f\) and all their derivatives decay exponentially on
the end, with constants depending on fixed \(N\).  Moreover
\begin{equation}
 t,Z=O_j(r^{2-n})\quad\hbox{for every }j,\qquad
 \Delta_gt=O(r^{-n-\beta}).
 \label{pde:end-decay}
\end{equation}
The uniform height bound and the radial height estimate of
Lemma~\ref{alg:trace-height} hold on the limiting solution.
\end{proposition}

\begin{proof}
First truncate at a coordinate sphere of radius \(R\), with
\(R\ge R_{\min}(N)\) as in the floor estimate.  For a trial
\(Z\in C^{1,\alpha}\) with zero outer trace solve the trace equation
by Lemma~\ref{pde:trial-trace}.  Freeze this pair in the divergence
equation, except for the gradient of the new \(Z\) in its principal
flux.  The resulting linear Dirichlet solve defines a map \(T\)
into \(C^{2,\alpha}\).  Indeed its principal coefficients are
\(C^{1,\alpha}\), its drift flux is \(C^{1,\alpha}\), and its
scalar source is \(C^\alpha\).  Thus \(T\) is a continuous compact
map on \(C^{1,\alpha}\).  Fixed points bootstrap to smoothness.

Use the homotopy described with \eqref{alg:system}, in reverse order:
first take \(K=0\) and scale the full scalar source from zero to one,
then scale \(K\) from zero to its specified value with the full source.
When \(K=0\), uniqueness of the trace equation gives \(f=0\) for
every trial.  At zero scalar scaling the map is identically zero.

For a homotopy parameter \(\lambda\), write \(f_\lambda[Z]\) for
the trace solve.  Apply degree on the relatively open moving domain
\[
 \left\{(\lambda,Z):
   \min t[Z,Df_\lambda[Z]]>-\varepsilon,\quad
   \|Z\|_{C^{1,\alpha}}<M_{N,R}\right\}.
\]
The dependence is continuous by Lemma~\ref{pde:trial-trace}.
Choose \(M_{N,R}\) larger than the fixed-point bound supplied by
Lemma~\ref{pde:coupled-estimates}.  There is no norm-boundary fixed
point.  There is no floor-boundary fixed point by the floor exclusion,
including its outer-sphere estimate and final \(K=f=0\) stage.
The a priori estimates also hold with a non-strict floor inequality
when taking a closure, before that exclusion is applied.
The map has relatively compact image on the bounded trial closure.
Leray--Schauder degree for a bounded relatively open subset of the
product with the parameter interval therefore applies
\cite[Sections~12--13]{LeraySchauder1934}.  Its degree at the initial
zero map is one, by excision around \(Z=0\); it stays one throughout
the homotopy.  This gives a smooth solution above the floor on each
sufficiently large truncation.

Lemma~\ref{pde:coupled-estimates} gives smooth local compactness
uniform in \(R\).  To retain the boundary condition at infinity,
we establish uniform end estimates before exhausting.
On the end \(K=0\), and \eqref{pde:normalized-trace} becomes
\[
 A_\chi:\nabla^2f=\frac{\eta\sqrt D}{l}f,\qquad
 \frac{\eta\sqrt D}{l}\ge\eta/e.
\]
For fixed \(N\), sufficiently small \(\kappa_N>0\), and a sufficiently
large fixed inner radius, the functions
\(C_Ne^{-\kappa_N(r-r_0)}\) are barriers for both signs of \(f\).
The positive zeroth-order coefficient dominates their radial second
derivatives; enlarge \(r_0\) to control the radial Hessian error.
They dominate at the inner sphere and at the outer zero boundary.
Consequently \(f\) decays exponentially, uniformly in \(R\).
Local and outer-face estimates give the same decay for every derivative.
The implicit relation then gives exponentially small \(t-Z\), with
all derivatives.

Up to exponentially decaying errors, the end equation is
\begin{equation}
 k\Delta_gZ+
 k\big[(k-1)+p(Z)-\delta_0s_{p(Z)}\big]|dZ|_g^2
 =\rho-\vartheta(N(Z+\varepsilon))C_N\rho_0,
 \qquad s_p=p+(k-1)/2.
 \label{pde:end-equation}
\end{equation}
Set \(a_N(z)=(k-1)+p(z)-\delta_0s_{p(z)}\), and choose
\(\Phi_N(0)=0\), \(\Phi_N'(0)=1\), with
\(\Phi_N''/\Phi_N'=a_N\).
On the bounded range of \(Z\), \(\Phi_N'\) has positive lower and
upper bounds depending on fixed \(N\).  Thus \eqref{pde:end-equation}
implies
\[
 \Delta_g\Phi_N(Z)=O_N(r^{-n-\beta})+O_N(e^{-\kappa_Nr}).
\]
Comparison with multiples of
\(r^{2-n}(1-r^{-\beta/2})\), on each truncation, gives
\(\Phi_N(Z)=O_N(r^{2-n})\), hence \(Z=O_N(r^{2-n})\).
These barriers have the required signed Laplacian for large \(r\);
their \(r^{-n-\beta/2}\) term dominates both the source and the
asymptotically flat metric errors.

Scaled Poisson \(W^{2,s}\) estimates first give weighted first-derivative
H\"older bounds; the source is a smooth function of \(Z\), the radius
weights and exponentially small terms.  Scaled Schauder and
differentiation then give \(Z=O_j(r^{2-n})\) for every \(j\).
The same follows for \(t\).  Its Laplacian has the improved bound in
\eqref{pde:end-decay}, since \(|dZ|^2=O(r^{2-2n})\) and
\(\beta<1\).  Exhaustion now supplies the asserted solution with
zero limit at infinity.  If the limit touched \(-\varepsilon\), that
point would be an interior minimum, and the identical floor calculation
of Lemma~\ref{alg:floor} would exclude it.  All height bounds pass to
the limit.
\end{proof}

\subsection{Mass, height separation, and the charged region}

\begin{proposition}[The mass cost of the filled construction]
\label{pde:mass-cost}
Let \(\widehat g=e^{2t}(g+l^2df^2)\) be obtained from
Proposition~\ref{pde:filled-solve}, and let \(E_*\) be the ADM energy
of the prepared metric \(g\).  Then
\begin{equation}
 \mathfrak F
 :=\lim_{r\to\infty}\int_{S_r}g(V,\nu)\,\dd A_g
 =\int_{\mathcal M_N}u\Xi\,\dd V_g,\qquad
 \mathcal E_{\mathrm{ADM}}(\widehat g)
 =E_*-\frac{\mathfrak F}{k\omega}.
 \label{pde:mass-flux}
\end{equation}
For fixed \(\varepsilon\),
\begin{equation}
 \mathfrak F\ge
 -\mathcal P_N(\ell+\ell^2/\eta)
 =-\mathcal P_N(\ell+\ell^{1/2})=-o_N(1).
 \label{pde:mass-error}
\end{equation}
The metric \(g_b=e^{2\varepsilon}\widehat g\), written in its
dilated asymptotic coordinates, consequently has
\begin{equation}
 \mathcal E_{\mathrm{ADM}}(g_b)
 \le e^{(n-2)\varepsilon}(E_*+o_N(1)).
 \label{pde:base-mass}
\end{equation}
It has all-order metric decay \(O_j(r^{2-n})\), a finite ADM energy,
and scalar curvature \(O(r^{-n-\beta})\).
\end{proposition}

\begin{proof}
All terms in \(u\Xi\) are integrable.  On the end this follows from
the exponential decay of \(f\),
\(\mathcal T=O(r^{2-2n})\), and the integrable radius weight
\(\rho_0=r^{-n-\beta}\); on the filled compact part it is ordinary
smooth integrability.  The divergence theorem on the manifold without
an inner boundary proves the flux limit in \eqref{pde:mass-flux}.
At infinity \(V=k\nabla t+o(r^{1-n})\) in the flux sense, since
\(u=1+O_N(r^{2-n})\) and all graph terms decay exponentially.
The metric perturbation \(2t\delta_{ij}\) contributes
\(-2k\,\partial_i t\) to the Euclidean ADM integrand.
The normalization \(1/(2k\omega)\) therefore gives the second
identity in \eqref{pde:mass-flux}; all mixed error fluxes vanish.

Only the penalty band can contribute negatively to the integral.
There \(m_0\ne0\) implies
\(-\varepsilon<t<-\varepsilon+1/N\), and hence
\(\ell\le l\le e\ell\) for sufficiently large \(N\).
Writing \(\sigma=|\nabla f|_g\), the definitions give
\begin{equation}
 u\le C(\ell+\ell^2\sigma),\qquad
 uv\le C\ell^2\sigma,\qquad
 ua_w\sigma\ge c\ell^2\sigma^2.
 \label{pde:penalty-band}
\end{equation}
For example the last expression equals
\(e^{(k-1)t}l^2\sigma^2\) exactly.
The negative penalty is bounded by
\[
 C\mathcal P_N\ell\rho_0
 +C\mathcal P_N\ell^2\sigma(\rho_0+\rho_1).
\]
Absorb its second term into the positive
\(\delta_0\eta ua_w\sigma\), using
\eqref{pde:penalty-band}.  The remaining negative part is bounded by
\[
 C\mathcal P_N\ell\rho_0
 +C\mathcal P_N\frac{\ell^2}{\eta}(\rho_0+\rho_1)^2.
\]
The integrals of \(\rho_0,\rho_0^2,\rho_1^2\) grow at most
polynomially in \(N\): the core volume is \(O(1+N)\), and on the end
\(4b_1>n\).  This proves \eqref{pde:mass-error}.  Notice that no
fixed-\(N\) Schauder constant enters this estimate.

A constant metric dilation by \(e^{2\varepsilon}\), followed by the
corresponding coordinate dilation, multiplies ADM energy by
\(e^{(n-2)\varepsilon}\), proving \eqref{pde:base-mass}.
All-order metric decay follows from
Proposition~\ref{pde:filled-solve} and the prepared end.
The conformal scalar law, the prepared scalar tail, and
\(\Delta_gt=O(r^{-n-\beta})\) give the stated scalar decay.
\end{proof}

\begin{proposition}[Separation of the original boundary by height]
\label{pde:height-separation}
There are \(b_2>0\) and \(N_2\), depending on the fixed prepared data
and \(\varepsilon\), such that the solutions of
Proposition~\ref{pde:filled-solve} satisfy, for \(N\ge N_2\),
\[
 h<-b_2\quad\hbox{on every plus component of }S,\qquad
 h>b_2\quad\hbox{on every minus component of }S.
\]
The constants are independent of \(N\).
For every fixed \(b>0\), the portion of the original exterior on which
\(|h|\ge b\) is contained in a fixed compact set, independent of \(N\).
\end{proposition}

\begin{proof}
The last assertion follows immediately from the uniform height tail
in Lemma~\ref{alg:trace-height}.  We prove the first assertion by a
trace supersolution on each long product and its prescribed collars.
Only the trace equation is used in this comparison, so no sign
condition in the distant cap is needed.

Treat a plus component.  Choose a small fixed \(\delta>0\) below the
strict negative leaf-expansion margin on the prescribed collars.
On the middle of its long product use height \(-\delta\).
There \(K=-L_0g\), with \(L_0\) the fixed large number used when
constructing the filling.  Towards the original boundary let the
slope increase from zero to a small positive fixed value \(a\).
It can be chosen as a smoothly truncated exponential: start with
size of order \(\eta\), increase exponentially, and stop smoothly
at \(a\), with
\[
 |(\hbox{slope})'|\le C(\hbox{slope}+\eta).
\]
Its length is \(O(1+|\log\eta|)=O(1+N)\), and its total height
increase is as small as desired by choosing \(a\) small.
The initial smoothing from slope zero has the same displayed bound.

For a height test \(H\), the Hessian multiplier in the trace equation
is
\[
 \frac{1}{\sqrt{(\eta/l)^2+|dH|^2}}.
\]
Thus the axial Hessian contribution along this ramp is uniformly
bounded, since \(l\le e\) and
\[
 \frac{\hbox{slope}+\eta}
      {\sqrt{(\eta/l)^2+\hbox{slope}^2}}\le C.
\]
The transverse contribution of the product metric is zero.
The trace of \(-L_0g\) is at most \(-kL_0\).  Taking \(L_0\)
large compared with the fixed ramp bound makes the ramp a strict
supersolution while its height is at least \(-\delta\).
This choice is part of the fixed filling geometry, before \(N\).

Continue at small positive slope through the fixed transitions to
\(S\), keeping the height there below \(-\delta/2\).
Then, on a fixed strict trapped collar on the original side,
increase the slope smoothly until the height exceeds the uniform
coarse height bound.  All these slopes are bounded below by \(a>0\);
the slopes and their derivatives on this fixed portion are fixed
finite numbers.  At a contact with the solution,
\(\eta/l\le\ell^{1/2}\), and the axial eigenvalue satisfies
\[
 \chi\le d
 =\frac{(\eta/l)^2}{(\eta/l)^2+|dH|^2}.
\]
Consequently the trace operator on the test converges uniformly,
on that fixed portion, to the expansion of its coordinate leaf,
\(H_{\mathrm{leaf}}+\operatorname{tr}_{\mathrm{leaf}}K\).
It is strictly negative by construction.  The error from the
axial second derivative tends to zero as \(N\to\infty\).
With the original choice of \(\delta\), this trace is below the
test height, including while the height increases to dominate the
coarse bound.

On the side of the long product facing the cap, mirror the ramp
and continue farther into product geometry at fixed slope, with
height increasing until it also dominates the coarse bound.
This adds only a fixed amount of length after the \(O(N)\) ramp.
It closes the comparison region without assuming anything about
the values of \(h\) in the cap, even when opposite-sign products
join there.

The test therefore dominates \(h\) on both ends of the comparison
region and is a strict supersolution within it.
At a positive interior maximum of \(h-H\), the two gradients agree.
The implicit \(t\), \(l\), and \(A_\chi\) consequently agree when
evaluating the test with the actual \(Z\), while the Hessian of
\(h-H\) is nonpositive.  The trace equation would give
\(h\le F(H)<H\), a contradiction.
Thus \(h\le H<-\delta/2\) on \(S\).
For a minus component apply the same argument to \(-h,-K\):
\(D,t,l,\chi\) are unchanged under the reversal, and the trace
equation reverses sign.  Taking the minimum over finitely many
components gives \(b_2\).  These barriers concern the final system;
the homotopy stages need not satisfy this comparison.
\end{proof}

\begin{proposition}[The charged low-height region]
\label{pde:charged-region}
Choose \(b_3>0\) sufficiently small relative to \(b_2\) and the
strict prepared matter margin.  For every sufficiently large \(N\),
put
\[
 \bar g=g+l^2df^2,\qquad
 g_b=e^{2(t+\varepsilon)}\bar g.
\]
On the original exterior \(g_b\ge g\) as quadratic forms, and
\[
 |dh|_{\bar g}\le\eta/l\le\ell^{1/2}.
\]
If \(\alpha=\iota_EdV_g\) is the retained electric flux form, then
on the original region \(\{|h|\le b_3\}\),
\begin{equation}
 \frac{R_{g_b}}2\ge c_n^2|\alpha|_{g_b}^2.
 \label{pde:charged-curvature}
\end{equation}
There is \(c_*>0\), independent of \(N\), such that on the original
band \(b_3/4\le|h|\le b_3\),
\begin{equation}
 e^{2(t+\varepsilon)}\frac{R_{g_b}}2
 \ge I_g(w)+c_*+(1-\delta_0)\mathcal T.
 \label{pde:retained-band}
\end{equation}
The band lies in a fixed compact portion of the original exterior.
Here \(I_g(w)\) is the electromagnetic expression defined with the
prepared data; it also permits the compactly supported magnetic form
used in the strict numerical extension.
\end{proposition}

\begin{proof}
The strict margin and the choice of \(\rho\) in
\eqref{alg:system} give
\(\mu_m-|J_m|-\rho>0\).  Since \(\tau\) is compactly supported,
we may choose \(b_3<b_2\), independent of \(N\), so that
\[
 \mu_m-|J_m|-\rho-b_3|\tau|
 \ge\tfrac12(\mu_m-|J_m|-\rho)>0
\]
everywhere on the original exterior.
The scalar identity \eqref{alg:scalar-identity}, the two equations,
and \(w(F)=\eta a_w|\nabla f|\) give
\begin{align*}
 e^{2t}\frac{R_{\widehat g}}2
 ={}&\mu+J(w)-\rho-h\tau+(1-\delta_0)\mathcal T\\
 &+(1-\delta_0)\eta a_w|\nabla f|
     +m_0C_N(\rho_0+v\rho_1).
\end{align*}
The DEC in the retained electromagnetic variables implies
\(\mu+J(w)\ge I_g(w)+\mu_m-|J_m|\).
All displayed terms on the second line are nonnegative.
This proves a positive matter remainder and the retained
\((1-\delta_0)\mathcal T\) throughout the low region.
On the stated band its lower bound is uniform: by
Proposition~\ref{pde:height-separation} the band is contained in a
fixed compact set.  This proves \eqref{pde:retained-band}, since
\(R_{g_b}=e^{-2\varepsilon}R_{\widehat g}\).

The inverse of \(\bar g\) is \(A_d\).  Hence
\[
 |dh|_{\bar g}^2
 =\eta^2d|\nabla f|_g^2
 =(\eta/l)^2v\le(\eta/l)^2.
\]
For the \((n-1)\)-form \(\alpha=\iota_EdV_g\), the volume factor
\(\sqrt D\) gives the exact norm identity
\begin{equation}
 |\alpha|_{\bar g}^2
 =d|E|_g^2+\langle w,E\rangle_g^2
 =|E|_g^2-|w\wedge E|_g^2.
 \label{pde:flux-form-norm}
\end{equation}
The electromagnetic square completion gives
\(I_g(w)\ge c_n^2|\alpha|_{\bar g}^2\).
Set \(z=t+\varepsilon>0\).  Scalar curvature scales by \(e^{-2z}\)
in the preceding inequality, whereas the squared norm of an
\((n-1)\)-form scales by \(e^{-2kz}\).  Since \(k\ge1\),
\[
 \frac{R_{g_b}}2
 \ge c_n^2e^{-2z}|\alpha|_{\bar g}^2
 \ge c_n^2e^{-2kz}|\alpha|_{\bar g}^2
 =c_n^2|\alpha|_{g_b}^2.
\]
Finally \(e^{2z}\ge1\) and \(\bar g\ge g\), proving metric dominance.
\end{proof}
